\documentclass[10pt, letterpaper]{article}
\usepackage[utf8]{inputenc}
\usepackage[margin=0.75in]{geometry}
\usepackage{enumitem}
\usepackage{hyperref}
\usepackage{titlesec}

% Formattazione delle sezioni
\titleformat{\section}{\large\bfseries\uppercase}{}{0em}{}[\titlerule]
\titlespacing{\section}{0pt}{1ex}{1ex}

\begin{document}

% Intestazione
\begin{center}
    {\LARGE \textbf{MAURIZIO DONNA}}\\[4pt]
    {\large SENIOR FPGA SYSTEM ARCHITECT}\\[4pt]
    PALO ALTO, CA | +1 650 304 9266 | \href{mailto:ing.maurizio.donna@gmail.com}{ing.maurizio.donna@gmail.com}\\
    \href{https://linkedin.com/in/mauriziodonnaing}{linkedin.com/in/mauriziodonnaing}
\end{center}

% Profilo
\section*{Professional Summary}
Senior FPGA System Architect with over 22 years of design expertise and a proven track record of architectural ownership across Silicon Valley and European markets. Possesses 18 years of experience as an authorized FPGA technical trainer and 7 years leading cross-functional engineering teams utilizing SAFe Agile and Scrum methodologies. Specializes in bridging complex hardware description languages with modern verification frameworks, custom PCB design, and automated CI/CD pipelines to deliver high-reliability systems for advanced research, telecommunications, and high-speed broadcasting.

% Competenze Tecniche
\section*{Technical Skills}
\begin{itemize}[leftmargin=*, nosep]
    \item \textbf{Architectures:} AMD/Xilinx (Kintex-7, Virtex-6/7, KU/VU+, Spartan), Lattice Semiconductor, Altera (Stratix II/V).
    \item \textbf{Languages:} VHDL, SystemVerilog, Verilog, C, Python, Bash, Tcl, SystemC.
    \item \textbf{Simulation \& Verification:} Cocotb, GHDL, OSVVM, ModelSim, QuestaSim, Active-HDL, ISIM.
    \item \textbf{Implementation \& Tools:} Vivado, Quartus II, ispLEVER, Synplify (Pro/DSP), XST, ChipScope, Matlab, Simulink.
    \item \textbf{Protocols \& Interfaces:} PCIe, GEth, SDH, OTU, AXI, SPI, I2C, ATSC.
    \item \textbf{DevOps \& Methodologies:} CI/CD (GitHub Actions, GitLab CI, Podman), Git, Mercurial, Jenkins, Jira, Linux, SAFe Agile, Scrum.
\end{itemize}

% Esperienza Professionale
\section*{Professional Experience}

\noindent \textbf{SLAC NATIONAL ACCELERATOR LABORATORY} | MENLO PARK, CA\\
\textit{FPGA Engineer, Accelerator Directorate - Electronics Engineering Division} | September 2022 - Present
\begin{itemize}[leftmargin=*, nosep]
    \item Direct the architectural vision for advanced accelerator instrumentation, establishing hardware-software co-design strategies leveraging AMD Kintex-7 and Lattice MachXO3LF FPGAs.
    \item Lead end-to-end custom PCB hardware architecture, defining strict 24V system requirements and integrating complex power monitoring and serial communication protocols to ensure system reliability.
    \item Spearheaded the modernization of the division's engineering workflows by establishing an automated CI/CD framework (GitHub Actions, GitLab CI, Podman), significantly reducing integration bottlenecks and setting new deployment standards.
    \item Develop robust, scalable testing and verification frameworks leveraging SystemVerilog, Cocotb, and GHDL.
    \item Mentor junior developers, coordinate project tasks, and implement best practices to ensure timely delivery of high-quality R\&D design.
\end{itemize}

\vspace{6pt}
\noindent \textbf{HITACHI ENERGY} | LUDVIKA, SWEDEN\\
\textit{FPGA Expert \& Scrum Master} | January 2020 - February 2022
\begin{itemize}[leftmargin=*, nosep]
    \item Owned the firmware lifecycle for critical MACH/VALVE control platforms, defining upgrade paths, resolving complex architectural bottlenecks, and ensuring high-availability performance.
    \item Spearheaded SAFe Agile transformations as Scrum Master, guiding cross-functional R\&D teams through complex development sprints and aligning firmware roadmaps with broader business objectives.
\end{itemize}

\vspace{6pt}
\noindent \textbf{SILICOM DENMARK A/S (FORMERLY FIBERBLAZE)} | SØBORG, DENMARK\\
\textit{Senior FPGA Architect \& Developer} | November 2017 - January 2020
\begin{itemize}[leftmargin=*, nosep]
    \item Led the architectural design and development of high-performance FPGA capture solutions targeting advanced Xilinx architectures, including Virtex-6, Virtex-7, and Kintex/Virtex UltraScale+.
    \item Executed critical SDAccel porting initiatives for custom hardware cards to optimize data center performance.
\end{itemize}

\vspace{6pt}
\noindent \textbf{EUROPEAN SPALLATION SOURCE ERIC} | LUND, SWEDEN\\
\textit{Engineer, Beam Diagnostic} | April 2014 - October 2017
\begin{itemize}[leftmargin=*, nosep]
    \item Designed the core FPGA architecture for Fast Beam Interlock instrumentation, deploying highly reliable Beam Current, Beam Position, and Beam Loss Monitors.
    \item Drove the adoption of a standardized, instrument-common PCIe framework on the MTCA.4 platform, creating a scalable, reusable architecture that unified hardware communication across the facility.
\end{itemize}

\vspace{6pt}
\noindent \textbf{INDEPENDENT CONSULTANT} | EUROPE\\
\textit{Freelance FPGA Expert \& IP Developer} | January 2008 - April 2014
\begin{itemize}[leftmargin=*, nosep]
    \item \textbf{Anritsu:} Designed a complete SDH LP for AU-3 and AU-4 with Ethernet statistics tracking, utilizing Altera Stratix V architectures.
    \item \textbf{STMicroelectronics:} Executed rigorous ASIC validation for specific intellectual properties using Altera Stratix II and Xilinx Virtex 5 FPGAs.
    \item \textbf{Durst:} Defined the architecture and implemented a complex communication board utilizing PCIe, DDR2, and Aurora protocols on a Virtex-5.
    \item \textbf{ABE:} Design porting from SPARTAN3 to SPARTAN 3DSP and improvement on a DVBT transceiver, implementation of a modulator based on ATSC standard on SPARTAN3DSP.
    \item \textbf{Additional Engagements:} Delivered hardware porting, proprietary algorithm implementation, and telecom standard (TIA/ATSC) modem designs for defense and telecom leaders including MBDA, Selex SI, and Lattice Italy.
\end{itemize}

\vspace{6pt}
\noindent \textbf{PIRELLI BROADBAND SOLUTIONS} | MILAN, ITALY\\
\textit{FPGA Specialist} | July 2004 - September 2006
\begin{itemize}[leftmargin=*, nosep]
    \item Implemented a metropolitan access system (Project City8) on a Lattice ORT42G5 FPGA, successfully multiplexing and regenerating GbE, FC, ESCON, and SDH protocols.
\end{itemize}

\vspace{6pt}
\noindent \textbf{BANCA SELLA} | ITALY\\
\textit{Project Manager, International Circuits} | April 2001 - April 2003
\begin{itemize}[leftmargin=*, nosep]
    \item Managed network connections to VISA and administered hardware/software for ATM, POS, and e-commerce authorization systems.
\end{itemize}

% Istruzione
\section*{Education}
\noindent \textbf{CEFRIEL - POLITECNICO DI MILANO}\\
\textit{Master Level II (EQF 8): Digital Electronics and EDA} | 2003-2004
\begin{itemize}[leftmargin=*, nosep]
    \item Thesis: Analysis, specification, and partial implementation of the physical layer of an IEEE 802.16a transceiver.
\end{itemize}

\vspace{6pt}
\noindent \textbf{POLITECNICO DI TORINO}\\
\textit{Master's Degree (Laurea) in Electronic Engineering} | 1992-2000
\begin{itemize}[leftmargin=*, nosep]
    \item Thesis: Utilizing FPGA and DSP to analyze SEMG during walking.
\end{itemize}

% Pubblicazioni e Formazione Specializzata
\section*{Publications \& Specialized Training}
\begin{itemize}[leftmargin=*, nosep]
    \item \textbf{Xilinx Authorized Training Provider (2005 - Present):} Mentored and upskilled enterprise engineering teams across Europe, serving as a designated technical authority on complex FPGA architectures, DSP implementation, and advanced HDL coding standards.
    \item \textbf{USPAS Certification (2017):} Fundamentals of Accelerator Physics and Technology with Simulations and Measurements Lab (sponsored by UC Davis).
    \item \textbf{Key Publications:} Co-authored research including "The ESS FPGA Framework and its Application on the ESS LLRF System," "FPGA based hybrid computing platform for ESS linac simulator," and "FW/SW framework for SRF cavity active resonance control".
\end{itemize}

\end{document}
